% article document
\documentclass[12pt]{article}

% document formatting packages
\usepackage[utf8]{inputenc}
\usepackage{graphicx}
\usepackage{caption}
\usepackage{subcaption}
\usepackage[margin=1in]{geometry}
\usepackage{setspace}
\usepackage{wrapfig}
\usepackage{textcomp, gensymb}
\usepackage{booktabs}
\usepackage{parskip}

% miscellaneous packages
\usepackage{lipsum}


\begin{document}

Abstract of (TITLE OF DISSERTATION), by (AUTHOR'S NAME), Ph.D., Brown University, May (YEAR IN WHICH DEGREE IS TO BE AWARDED)

\vspace{2cm}
\doublespacing
From the dissertation guidelines (https://www.brown.edu/academics/gradschool/dissertation-guidelines):

The dissertation must be accompanied by an abstract. The abstract should, in a concise manner, present the problem of the dissertation, discuss the materials and procedure or methods used, and state the results or conclusions. Mathematical formulas, diagrams, and other illustrative materials should be avoided. The abstract should not be part of the dissertation itself nor should it be included in the table of contents. The abstract should be presented in two unnumbered loose copies. The abstract should be prepared carefully since it will be published without editing or revision. The abstract should be double-spaced and may not exceed 350 words (maximum 2,450 typewritten characters — including spaces and punctuation — about 70 characters per line with a maximum of 35 lines). 

\end{document}